% Part: model-theory
% Chapter: models-of-arithmetic
% Section: non-standard-models

\documentclass[../../../include/open-logic-section]{subfiles}

\begin{document}

\olfileid{mod}{mar}{nst}
\section{अमानक प्रतिमाने}

\begin{explain}
$\Lang{L_A}$ साठीची एखादी !!a{structure}~$\Struct{N}$ शी समरूप असेल,
तर तिला मानक म्हणतो. एखादी !!a{structure}~$\Struct{N}$ शी समरूप नसेल,
तर तिला अमानक म्हणतो.
\end{explain}

\begin{defn}
$\Lang{L_A}$ साठीची !!{structure}~$\Struct{M}$ जर~$\Struct{N}$ शी समरूप
नसेल, तर ती \emph{अमानक} आहे. प्रांतातील जे !!{element}s
$x \in \Domain{M}$ कोणत्यातरी $n \in \Nat$ साठी
$\Value{\num{n}}{M}$ बरोबर आहेत, त्यांना ($\Struct{M}$ मधील)
\emph{मानक संख्या} म्हणतात; आणि जे तसे नाहीत त्यांना \emph{अमानक संख्या}
म्हणतात.
\end{defn}

\begin{explain}
\olref[stm]{prop:standard-domain} नुसार, $\Lang{L_A}$ साठीच्या कोणत्याही
मानक !!{structure} मध्ये फक्त मानक !!{element}s असतात. त्यामुळे अमानक
!!{structure} मध्ये किमान एक अमानक घटक असलाच पाहिजे. किंबहुना, अमानक
!!{element} अस्तित्वात असणे हे ती !!{structure} अमानक असल्याची हमी देते.
\end{explain}

\begin{prop}
$\Lang{L_A}$ साठीच्या !!a{structure}~$\Struct{M}$ मध्ये अमानक संख्या
असेल, तर $\Struct{M}$ अमानक आहे.
\end{prop}

\begin{proof}
तसे नाही असे समजा; म्हणजे $\Struct{M}$ मानक असूनही तिच्यात अमानक संख्या
$x$ आहे असे समजा. $g\colon \Nat \to \Domain{M}$ हे समरूपण असू द्या.
$n$ वरील विगमनाने $g(\Value{\num{n}}{N}) = \Value{\num{n}}{M}$ हे सहज
दिसते. दुसऱ्या शब्दांत, $g$ हे~$\Struct{N}$ मधील मानक संख्यांना
$\Struct{M}$ मधील मानक संख्यांवर पाठवते. $\Struct{M}$ मध्ये अमानक संख्या
असेल, तर गृहीतकाच्या विरुद्ध $g$ हे !!{surjective} असू शकत नाही.
\end{proof}

\begin{prob}
$\Th{Q}$ मध्ये पुढील स्वयंसिद्धके आहेत हे आठवा:
\begin{align*}
& \lforall[x][\lforall[y][(\eq[x'][y'] \lif \eq[x][y])]] \tag{$!Q_1$}\\
& \lforall[x][\eq/[\Obj 0][x']] \tag{$!Q_2$}\\
& \lforall[x][(\eq[x][\Obj 0] \lor \lexists[y][\eq[x][y']])] \tag{$!Q_3$}
\end{align*}
पुढील अटी पूर्ण करणाऱ्या !!{structure}s~$\Struct{M_1}$, $\Struct{M_2}$,
$\Struct{M_3}$ द्या:
\begin{enumerate}
\item $\Sat{M_1}{!Q_1}$, $\Sat{M_1}{!Q_2}$, $\Sat/{M_1}{!Q_3}$;
\item $\Sat{M_2}{!Q_1}$, $\Sat/{M_2}{!Q_2}$, $\Sat{M_2}{!Q_3}$; आणि
\item $\Sat/{M_3}{!Q_1}$, $\Sat{M_3}{!Q_2}$, $\Sat{M_3}{!Q_3}$.
\end{enumerate}
अर्थात, प्रत्येकासाठी फक्त~$\Assign{\Obj{0}}{M_i}$ आणि
$\Assign{\prime}{M_i}$ निश्चित करायची आहेत.
\end{prob}

\begin{explain}
$\Lang{L_A}$ साठी अमानक !!{structure}s निश्चित करणे पुरेसे सोपे आहे.
उदाहरणार्थ, !!{domain}~$\Int$ असलेली रचना घ्या आणि सर्व अतार्किक चिन्हांचे
नेहमीप्रमाणे अर्थनिर्धारण करा. ऋण संख्या कोणत्याही~$n$ साठी $\num{n}$ ची
मूल्ये नसल्यामुळे ही रचना अमानक आहे. अर्थात, $\Struct{N}$ प्रमाणेच ती
वाक्ये सत्य करते या अर्थाने ती अंकगणिताचे \emph{प्रतिमान} नसेल. उदाहरणार्थ,
$\lforall[x][\eq/[x'][\Obj{0}]]$ हे असत्य आहे. तरीसुद्धा, संहतता प्रमेय
वापरून अंकगणिताची अमानक प्रतिमाने अस्तित्वात आहेत हे सहज सिद्ध करता येते.
\end{explain}

\begin{prop}
$\Th{TA} = \Setabs{!A}{\Sat{N}{!A}}$ ही~$\Struct{N}$ ची उपपत्ती असू द्या.
$\Th{TA}$ ला !!a{enumerable} अमानक प्रतिमान आहे.
\end{prop}

\begin{proof}
$\Lang{L_A}$ मध्ये नवीन !!{constant}~$c$ घालून तिचा विस्तार करा आणि
पुढील !!{sentence}s चा संच विचारात घ्या:
\[
\Gamma = \Th{TA} \cup \{\eq/[c][\num{0}], \eq/[c][\num{1}],
\eq/[c][\num{2}], \dots\}
\]
$\Gamma$ च्या कोणत्याही प्रतिमान~$\Struct{M^c}$ मध्ये
$x = \Assign{c}{M^c}$ असा अमानक !!a{element} असेल, कारण प्रत्येक
$n \in \Nat$ साठी $x \neq \Value{\num{n}}{M}$. तसेच
$\Th{TA} \subseteq \Gamma$ असल्यामुळे अर्थातच $\Sat{M^c}{\Th{TA}}$.
फक्त~$c$ विसरून $\Struct{M^c}$ पासून $\Lang{L_A}$ साठीची
!!a{structure}~$\Struct{M}$ बनवली, तरी तिच्या प्रांतात तो अमानक~$x$ राहतो
आणि $\Sat{M}{\Th{TA}}$ सुद्धा. $\Th{TA}$ मध्ये $c$ येत नसल्यामुळे या
दुसऱ्या दाव्याची हमी मिळते. म्हणून $\Gamma$ ला प्रतिमान आहे एवढे दाखवणे
पुरेसे आहे.
%READERNOTE{OLFOL-031 — गोठवलेल्या इंग्रजीतील पहिल्या सूत्रात c चे अर्थनिर्धारण M मध्ये घेतले आहे; पण c हे विस्तारित भाषेचे चिन्ह असल्यामुळे L_A-रचना M मध्ये ते परिभाषितच नाही. मराठी सूत्रात ते योग्य विस्तारित रचना M^c मध्ये घेतले आहे.}

$\Gamma$ ला प्रतिमान आहे हे दाखवण्यासाठी संहतता प्रमेय वापरू. $\Gamma$ चा
प्रत्येक सांत उपसंच पूर्ततायोग्य असेल, तर $\Gamma$ सुद्धा पूर्ततायोग्य आहे.
कोणताही सांत उपसंच $\Gamma_0 \subseteq \Gamma$ विचारात घ्या.
$\Gamma_0$ मध्ये~$\Th{TA}$ मधील काही !!{sentence}s आणि
$\eq/[c][\num{n}]$ या रूपातील सांत वाक्ये येतात. या दुसऱ्या रूपातील एकही
वाक्य नसेल, तर c ला शून्य मूल्य देणारा N चा विस्तार त्या उपसंचाला थेट
पूर्ण करतो; म्हणून पुढे किमान एक असे वाक्य आहे असे समजा.
%READERNOTE{OLFOL-033 — गोठवलेल्या इंग्रजी पुराव्यात मनमाना सांत उपसंच घेतल्यानंतर लगेच सर्वांत मोठा k निवडला आहे; पण त्या उपसंचात c-विषयक असमता एकही नसेल, तर असा k नसतो. मराठी पुराव्यात तो रिकामा प्रसंग स्वतंत्रपणे हाताळला आहे.}
$\eq/[c][\num{k}] \in \Gamma_0$ असा $k$ हा सर्वांत मोठा अंक समजा.
$\Assign{c}{N_k} = k+1$ हे अर्थनिर्धारण समाविष्ट करण्यासाठी~$\Struct{N}$
चा विस्तार करून $\Struct{N_k}$ ची व्याख्या करा. $\Sat{N_k}{\Gamma_0}$:
जर $!A \in \Th{TA}$, तर $\Struct{N_k}$ ही~$c$ वगळता प्रत्येक बाबतीत
$\Struct{N}$ सारखीच असल्यामुळे आणि~$!A$ मध्ये $c$ येत नसल्यामुळे
$\Sat{N_k}{!A}$. तसेच $n \le k$ आणि $\Value{c}{N_k} = k+1$ असल्यामुळे
$\Sat{N_k}{\eq/[c][\num{n}]}$. म्हणून $\Gamma$ चा प्रत्येक सांत उपसंच
पूर्ततायोग्य आहे आणि संहततेनुसार संपूर्ण संचाला प्रतिमान आहे. अधोगामी
लोव्हेनहाइम--स्कोलेम प्रमेयानुसार हे प्रतिमान गणनीय घेता येते.
%READERNOTE{OLFOL-032 — विधानात गणनीय अमानक प्रतिमानाचा दावा आहे; गोठवलेल्या इंग्रजी पुराव्यात संहततेने प्रतिमान मिळवल्यानंतर त्याची गणनीयता सिद्ध केलेली नाही. उपलब्ध पूर्वप्रमेयानुसार अधोगामी लोव्हेनहाइम--स्कोलेम प्रमेय लावून ती पायरी मराठीत स्पष्ट केली आहे.}
\end{proof}

\end{document}
